% AY2017 制御工学 第1問〔I〕（転記）
% 出典: 03_制御工学/original/03_制御工学/2017/2017se.pdf
% 要約: 二次振動系（質量 m, 粘性係数 d, バネ定数 k, 力 f(t), 変位 x(t)）。
%       運動方程式・伝達関数、ωn/ζ=0 設定からの d,k と自由応答、指数入力応答。

\section*{第1問〔I〕}

下図の二次振動系に関する次の問い (1)〜(3) に答えよ. なお, 質量を $m$, 粘性係数を $d$,
ばね定数を $k$ とし, 時間 $t$ において, $f(t)$ は系に加える力, $x(t)$ は質量の中心位置を表す.

\begin{center}
\begin{tikzpicture}[>=Latex]
  \draw[thick] (0,-0.2) -- (0,2.2);
  \foreach \y in {-0.1,0.1,...,2.05} \draw[thick] (0,\y) -- (-0.25,\y-0.25);
  \draw[thick,decorate,decoration={zigzag,segment length=8,amplitude=5}] (0,1.6) -- (3,1.6);
  \node at (1.5,2.05) {$k$};
  \draw[thick] (0,0.45) -- (1.15,0.45);
  \draw[thick] (1.15,0.2) rectangle (1.85,0.7);
  \draw[thick] (1.55,0.45) -- (3,0.45);
  \node at (1.5,-0.05) {$d$};
  \draw[thick,fill=gray!12] (3,0.1) rectangle (3.9,2.0);
  \node at (3.45,1.05) {$m$};
  \draw[->,thick] (3.9,1.05) -- (5.0,1.05);
  \node at (5.25,1.05) {$f(t)$};
  \draw[->,thick] (3.45,2.35) -- (4.3,2.35);
  \node at (3.9,2.65) {$x(t)$};
\end{tikzpicture}
\end{center}

\begin{enumerate}[label=(\arabic*)]
  \item 系の運動方程式を求めたうえで, 入力を $f(t)$, 出力を $x(t)$ として伝達関数を求めよ.

  \item 質量を $m=1$, 系の固有振動数を $\omega_n=1$, 減衰係数を $\zeta=0$ と設定する場合,
        $d$ ならびに $k$ を求めよ. また, この系に対して入力を $f(t)=0$, 初期値を $x(0)=1$,
        $\dot{x}(0)=0$ と設定したときの応答 $x(t)$ を求めたうえで, $x(t)$ の概形を示せ.
        なお, 横軸を時間 $t$, 縦軸を $x(t)$ とすること.

  \item 問い (2) のシステムに対して, 入力 $f(t)=2e^{-t}$ を印加した（ただし, $t<0$ のとき
        $f(t)=0$）. $x(t)$ を求めたうえで, 時間が充分に経過した時の振幅の最大値を求めよ.
\end{enumerate}
